
MeerTime is timing $\NMeerTimePSRs$ pulsars.
Using observations to date, an anticipated $\sigma_I$ and integration time (for a single pulsar observation) required to achieve it can be estimated for each pulsar.
The value of $\sigma_I$ scales with observation $T_{{\rm obs},I}$ time like 
\begin{equation}
\sigma_I = \delta_I / \sqrt{T_{{\rm obs},I}}\,,
\label{eqn:TobsToSig}
\end{equation}
where $\delta_I$ is a constant of proportionality which depends on the pulsar. 
Throughout this work we make the simplifying assumption that $\delta_I$ is a constant for each pulsar throughout the PTA lifespan \han{[maybe add some discussion here about this?]}

\subsection{Spin noise and jitter noise}
MeerTime observations are not currently long enough to estimate spin noise in every pulsar.
However spin noise measurements already exist for $13$ of the pulsars from PPTA and NANOGrav analysis. 
We take $\ARed$ and $\gammaRed$ values from ~\cite{GoncharovEtAl:2021} (sixth and seventh column in Table 1) and ~\cite{NANOGrav12p5yr:2021} (tenth and eleventh column in Table 5). 
Where a value is listed in both papers, we take the~\cite{GoncharovEtAl:2021} value. 
Table~\ref{tab:redNoise} shows the red noise parameters for $13$ of the $\NMeerTimePSRs$ pulsars. 
Otherwise, we assume zero spin noise. 

\begin{table}
\begin{center}
\begin{tabular}{llll}
\hline \hline
Pulsar     & $\log_{10}(\ARed)$  & $\gammaRed$ & Reference \\
\hline
J0030+0451 & $-15.72$ & $5.3$   &  \cite{NANOGrav12p5yr:2021}\\
J0437-4715 & $-14.56$ & $2.99$  &  \cite{GoncharovEtAl:2021}\\
J0613-0200 & $-14.26$ & $4.17$  &  \cite{GoncharovEtAl:2021}\\
J0711-6830 & $-13.04$ & $1.09$  &  \cite{GoncharovEtAl:2021}\\
J1024-0719 & $-14.62$ & $6.39$  &  \cite{GoncharovEtAl:2021}\\
J1600-3053 & $-14.34$ & $3.81$  &  \cite{GoncharovEtAl:2021}\\
J1643-1224 & $-12.85$ & $0.98$  &  \cite{GoncharovEtAl:2021}\\
J1713+0747 & $-15.16$ & $1.6$   &  \cite{NANOGrav12p5yr:2021}\\
J1744-1134 & $-14.16$ & $1.7$   &  \cite{NANOGrav12p5yr:2021}\\
J1747-4036 & $-13.78$ & $3.8$   &  \cite{NANOGrav12p5yr:2021}\\
J1909-3744 & $-14.74$ & $4.05$  &  \cite{GoncharovEtAl:2021}\\
J2145-0750 & $-13.86$ & $1.6$   &  \cite{NANOGrav12p5yr:2021}\\
J2317+1439 & $-16.50$ & $6.0$   &  \cite{NANOGrav12p5yr:2021}\\
\hline\hline
\end{tabular}
\end{center}
\caption{\label{tab:redNoise}
Spin noise parameters for $13$ pulsars. 
The second and third columns show the log of the spin noise amplitude and the slope. 
}
\end{table}

Some of the pulsars have jitter noise estimates. 
Table~\ref{tab:jitter} shows jitter noise measurements for $13$ MeerTime pulsars from~\cite{ParthasarathyEtAl:2021}.
We assume the rest of the pulars have zero jitter noise.


\begin{table}
\begin{center}
\begin{tabular}{lll}
\hline \hline
Pulsar     & $\sigJitter$   & Reference \\
\hline
J0125-2327 & $48  \pm 13$   & ~\cite{ParthasarathyEtAl:2021}\\
J0437-4715 & $50  \pm 10$   & ~\cite{ParthasarathyEtAl:2021}\\
J0636-3044 & $100 \pm 30$   & ~\cite{ParthasarathyEtAl:2021}\\
J0711-6830 & $60  \pm 20$   & ~\cite{ParthasarathyEtAl:2021}\\
J1022+1001 & $120 \pm 20$   & ~\cite{ParthasarathyEtAl:2021}\\
J1045-4509 & $130 \pm 75$   & ~\cite{ParthasarathyEtAl:2021}\\
J1603-7202 & $180 \pm 40$   & ~\cite{ParthasarathyEtAl:2021}\\
J1730-2304 & $80  \pm 45$   & ~\cite{ParthasarathyEtAl:2021}\\
J1744-1134 & $30  \pm  6$   & ~\cite{ParthasarathyEtAl:2021}\\
J1757-5322 & $130 \pm 45$   & ~\cite{ParthasarathyEtAl:2021}\\
J1909-3744 & $9   \pm  3$   & ~\cite{ParthasarathyEtAl:2021}\\
J2145-0750 & $200 \pm 20$   & ~\cite{ParthasarathyEtAl:2021}\\
J2241-5236 & $3.8 \pm  0.8$ & ~\cite{ParthasarathyEtAl:2021}\\
\hline\hline
\end{tabular}
\end{center}
\caption{\label{tab:jitter}
Jitter noise measurements for $13$ pulsars. 
}
\end{table}

% pulsar reference - jitter, white nooise, red noies, position

\begin{comment}
% value from NANOGRav table is An
% An  = 0.006 \mu s yr^{1/2} 
% An -> 0.006 * 3.17E-14 yr^{3/2}
\begin{table}
\begin{tabular}{llll}
\hline\hline
Pulsar     & $\ARed$ & $\gammaRed$ & Ref.\\
\hline
J0030+0451 & $-15.72$ & $5.3$   &  \cite{NANOGrav12p5yr:2021}\\
J0437-4715 & $-14.56$ & $2.99$  &  \cite{GoncharovEtAl:2021}\\
J0613-0200 & $-14.26$ & $4.17$  &  \cite{GoncharovEtAl:2021}\\
J0711-6830 & $-13.04$ & $1.09$  &  \cite{GoncharovEtAl:2021}\\
J1024-0719 & $-14.62$ & $6.39$  &  \cite{GoncharovEtAl:2021}\\
J1600-3053 & $-14.34$ & $3.81$  &  \cite{GoncharovEtAl:2021}\\
J1643-1224 & $-12.85$ & $0.98$  &  \cite{GoncharovEtAl:2021}\\
J1713+0747 & $-15.16$ & $1.6$   &  \cite{NANOGrav12p5yr:2021}\\
J1744-1134 & $-14.16$ & $1.7$   &  \cite{NANOGrav12p5yr:2021}\\
J1747-4036 & $-13.78$ & $3.8$   &  \cite{NANOGrav12p5yr:2021}\\
J1909-3744 & $-14.74$ & $4.05$  &  \cite{GoncharovEtAl:2021}\\
J2145-0750 & $-13.86$ & $1.6$   &  \cite{NANOGrav12p5yr:2021}\\
J2317+1439 & $-16.50$ & $6.0$   &  \cite{NANOGrav12p5yr:2021}\\
\hline
\end{tabular}
\caption{\label{tab:redNoise}
Spin noise parameters for $13$ pulsars. 
[Notes: Values from Ref.~\cite{GoncharovEtAl:2021} are taken directly form their Table 1. 
%Values from Ref.~\cite{NANOGrav12p5yr:2021} are converted from units of ${\rm \mu s \,yr}^{1/2}$ to $\log_{10} \ARed[{\rm yr}^{3/2}]$.]
}
\end{table}
\end{comment}

\begin{comment}
\textbf{estimating spin noise for other pulsars? }
Spin noise parameters for the remaining pulsars can be estimated using the scaling relation formalism in Shannon and Cordes. 
The spin noise RMS value is given by 
\begin{equation}
\sigma_{\rm SN} = C_2 \mu^{\alpha} | {\dot \mu}_{-15} |^{\beta} t_{\rm yr}^{\gamma}
\end{equation}
where $\mu$......
In this work we use the central fit values for the ${\rm MSP}_{\rm 10,PPTA}$+${\rm NANO}$ ($\{ \ln C_2, \alpha, \beta, \gamma \} = \{-1.2,-0.9,0.8,2.3\}$).
\end{comment}




\subsection{Current observation strategy}
MeerTime pulsars are observed with a cadence of once every two weeks ($c=26\,{\rm yr}^{-1})$. 
The total observation time used in each $2$ week observing time is $\TwoWkObsTimeS\,{\rm s} = \TwoWkObsTimeHrs\,{\rm hrs}$, which is divided between the MeerTime pulsars. 
The targeted precision for each pulsar is $\sigma_I = 1 {\rm \mu\, s}$. 

Using the current observation strategy and assuming a GWB signal with $A = 2\times 10^{-15}$ and $\alpha = -2/3$, we can estimate the predicted average SNR for MeerTime from Eqn.~\ref{eqn:snrPIPJPg}. 
After a timespan of $T=10\,{\rm yrs}$, the average SNR for the current obseravation stategy is $12.3$.
Here, we consider whether, or not, the average SNR can be improved by a re-allocation of observing time between pulsars. 




